\section{D}

% dauern (to last / take time)
\begin{verbentry}{
  style = {default},
  toc = {dauern (to last / take time)},
  name = {dauern},
  english = {to last / take time},
  type = {regular},
  auxiliary = {haben}
}
 
    
     \vspace{5pt}

    \hrule
    
    \vspace{5pt}

  \begin{tabular}{rl}
     \multicolumn{2}{l}{\textbf{PRÄSENS} }\\
    ich & \textbf{dauere}\\
    du & \textbf{dauerst}\\
    er/sie/es & \textbf{dauert}\\
    wir & \textbf{dauern}\\
    ihr & \textbf{dauert}\\
    sie/Sie & \textbf{dauern}\\
  \end{tabular}
  \hfill
     \begin{tabular}{rl}
    \multicolumn{2}{l}{\textbf{PRÄTERITUM} }\\
    ich & \textbf{dauerte}\\
    du & \textbf{dauertest}\\
    er/sie/es & \textbf{dauerte}\\
    wir & \textbf{dauerten}\\
    ihr & \textbf{dauertet}\\
    sie/Sie & \textbf{dauerten}\\
  \end{tabular}
  \hfill
  \begin{tabular}{rl}
     \multicolumn{2}{l}{\textbf{IMPERATIV} }\\
   & \textbf{dauere (du)!}\\
   & \textbf{dauern wir!}\\
   & \textbf{dauert (ihr)!}\\
   & \textbf{dauern Sie!}\\
   & \vspace{10pt}
  \end{tabular}

    \vspace{10pt}

 \begin{tabular}{rl}
     \multicolumn{2}{l}{\textbf{PERFEKT}}\\
   & haben + \textbf{gedauert}\\
  \end{tabular}
\hfill
   \begin{tabular}{rl}
    \multicolumn{2}{l}{\textbf{FUTURE I} }\\
   & werden + \textbf{dauern}\\
  \end{tabular}
\hfill
   \begin{tabular}{rl}
      \multicolumn{2}{l}{\textbf{PLUSQUAMPERFEKT}}\\
   & hatten + \textbf{gedauert}\\
  \end{tabular}

   \vspace{10pt}

   \begin{tabular}{lll}
      \multicolumn{3}{l}{\textbf{MIT PRÄPOSITIONEN}}\\
& \textbf{dauern bis} + Akk & (to last until)\\
& \textbf{dauern über} + Akk & (to last over / longer than)\\
   \end{tabular}
  \hfill
   \begin{tabular}{lll}
      \multicolumn{3}{l}{\textbf{TRENNBARE VERBEN}}\\
 & \textbf{an$|$dauern} & (to continue / persist)\\
\vspace{10pt}
  \end{tabular}

\vspace{-5pt}
\end{verbentry}

% demonstrieren (to demonstrate / protest)
\begin{verbentry}{
  style = {default},
  toc = {demonstrieren (to demonstrate / protest)},
  name = {demonstrieren},
  english = {to demonstrate / protest},
  type = {regular},
  auxiliary = {haben}
}


\vspace{5pt}
\hrule
\vspace{5pt}

\begin{tabular}{rl}
\multicolumn{2}{l}{\textbf{PRÄSENS}}\\
ich & \textbf{demonstriere}\\
du & \textbf{demonstrierst}\\
er/sie/es & \textbf{demonstriert}\\
wir & \textbf{demonstrieren}\\
ihr & \textbf{demonstriert}\\
sie/Sie & \textbf{demonstrieren}\\
\end{tabular}
\hfill
\begin{tabular}{rl}
\multicolumn{2}{l}{\textbf{PRÄTERITUM}}\\
ich & \textbf{demonstrierte}\\
du & \textbf{demonstriertest}\\
er/sie/es & \textbf{demonstrierte}\\
wir & \textbf{demonstrierten}\\
ihr & \textbf{demonstriertet}\\
sie/Sie & \textbf{demonstrierten}\\
\end{tabular}
\hfill
\begin{tabular}{rl}
\multicolumn{2}{l}{\textbf{IMPERATIV}}\\
& \textbf{demonstriere (du)!}\\
& \textbf{demonstrieren wir!}\\
& \textbf{demonstriert (ihr)!}\\
& \textbf{demonstrieren Sie!}\\
\end{tabular}

\vspace{10pt}

\begin{tabular}{rl}
\multicolumn{2}{l}{\textbf{PERFEKT}}\\
& haben + \textbf{demonstriert}\\
\end{tabular}
\hfill
\begin{tabular}{rl}
\multicolumn{2}{l}{\textbf{FUTURE I}}\\
& werden + \textbf{demonstrieren}\\
\end{tabular}
\hfill
\begin{tabular}{rl}
\multicolumn{2}{l}{\textbf{PLUSQUAMPERFEKT}}\\
& hatten + \textbf{demonstriert}\\
\end{tabular}

\vspace{10pt}

\begin{tabular}{lll}
\multicolumn{3}{l}{\textbf{MIT PRÄPOSITIONEN}}\\
& \textbf{demonstrieren gegen} + Akk & (to demonstrate against)\\
& \textbf{demonstrieren für} + Akk & (to demonstrate for)\\
\end{tabular}
\hfill
\begin{tabular}{lll}
\multicolumn{3}{l}{\textbf{TRENNBARE VERBEN}}\\
& \textbf{---} & \\
\end{tabular}

\vspace{-5pt}
\end{verbentry}

% denken (to think)
\begin{verbentry}{
  style = {acc},
  toc = {denken (to think)},
  name = {denken},
  english = {to think},
  type = {mixed, + Akkusativ},
  auxiliary = {haben}
}
 
    
     \vspace{5pt}

    \hrule
    
    \vspace{5pt}

  \begin{tabular}{rl}
     \multicolumn{2}{l}{\textbf{PRÄSENS} }\\
    ich & \textbf{denke}\\
    du & \textbf{denkst}\\
    er/sie/es & \textbf{denkt}\\
    wir & \textbf{denken}\\
    ihr & \textbf{denkt}\\
    sie/Sie & \textbf{denken}\\
  \end{tabular}
  \hfill
     \begin{tabular}{rl}
    \multicolumn{2}{l}{\textbf{PRÄTERITUM} }\\
    ich & \textbf{dachte}\\
    du & \textbf{dachtest}\\
    er/sie/es & \textbf{dachte}\\
    wir & \textbf{dachten}\\
    ihr & \textbf{dachtet}\\
    sie/Sie & \textbf{dachten}\\
  \end{tabular}
  \hfill
  \begin{tabular}{rl}
     \multicolumn{2}{l}{\textbf{IMPERATIV} }\\
   & \textbf{denk(e) (du)!}\\
   & \textbf{denken wir!}\\
   & \textbf{denkt (ihr)!}\\
   & \textbf{denken Sie!}\\
   & \vspace{10pt}
  \end{tabular}

    \vspace{10pt}

 \begin{tabular}{rl}
     \multicolumn{2}{l}{\textbf{PERFEKT}}\\
  &  haben + \textbf{gedacht}\\
  \end{tabular}
\hfill
   \begin{tabular}{rl}
   
    \multicolumn{2}{l}{\textbf{FUTURE I} }\\
  &  werden + \textbf{denken}\\
 
  \end{tabular}
\hfill
   \begin{tabular}{rl}
      \multicolumn{2}{l}{\textbf{PLUSQUAMPERFEKT}}\\
  &  hatten + \textbf{gedacht}\\
  \end{tabular}
  
  \vspace{10pt}

\begin{tabular}{lll}
\multicolumn{3}{l}{\textbf{MIT PRÄPOSITIONEN}}\\
& \textbf{---} & \\
\end{tabular}
\hfill
\begin{tabular}{lll}
\multicolumn{3}{l}{\textbf{TRENNBARE VERBEN}}\\
& \textbf{---} & \\
\end{tabular}

  \vspace{-5pt}
\end{verbentry}

% diskutieren (to discuss)
\begin{verbentry}{
  style = {acc},
  toc = {diskutieren (to discuss)},
  name = {diskutieren},
  english = {to discuss},
  type = {irregular, + Akkusativ},
  auxiliary = {haben}
}
 
    
     \vspace{5pt}

    \hrule
    
    \vspace{5pt}

  \begin{tabular}{rl}
     \multicolumn{2}{l}{\textbf{PRÄSENS} }\\
    ich & \textbf{diskutiere}\\
    du & \textbf{diskutierst}\\
    er/sie/es & \textbf{diskutiert}\\
    wir & \textbf{diskutieren}\\
    ihr & \textbf{diskutiert}\\
    sie/Sie & \textbf{diskutieren}\\
  \end{tabular}
  \hfill
     \begin{tabular}{rl}
    \multicolumn{2}{l}{\textbf{PRÄTERITUM} }\\
    ich & \textbf{diskutierte}\\
    du & \textbf{diskutiertest}\\
    er/sie/es & \textbf{diskutierte}\\
    wir & \textbf{diskutierten}\\
    ihr & \textbf{diskutiertet}\\
    sie/Sie & \textbf{diskutierten}\\
  \end{tabular}
  \hfill
  \begin{tabular}{rl}
     \multicolumn{2}{l}{\textbf{IMPERATIV} }\\
   & \textbf{diskutier(e) (du)!}\\
   & \textbf{diskutieren wir!}\\
   & \textbf{diskutiert (ihr)!}\\
   & \textbf{diskutieren Sie!}\\
   & \vspace{10pt}
  \end{tabular}

    \vspace{10pt}

 \begin{tabular}{rl}
     \multicolumn{2}{l}{\textbf{PERFEKT}}\\
  &  haben + \textbf{diskutiert}\\
  \end{tabular}
\hfill
   \begin{tabular}{rl}
   
    \multicolumn{2}{l}{\textbf{FUTURE I} }\\
  &  werden + \textbf{diskutieren}\\
 
  \end{tabular}
\hfill
   \begin{tabular}{rl}
      \multicolumn{2}{l}{\textbf{PLUSQUAMPERFEKT}}\\
  &  hatten + \textbf{diskutiert}\\
  \end{tabular}
  
  \vspace{10pt}

\begin{tabular}{lll}
\multicolumn{3}{l}{\textbf{MIT PRÄPOSITIONEN}}\\
& \textbf{---} & \\
\end{tabular}
\hfill
\begin{tabular}{lll}
\multicolumn{3}{l}{\textbf{TRENNBARE VERBEN}}\\
& \textbf{---} & \\
\end{tabular}

  \vspace{-5pt}
\end{verbentry}

% drehen (to turn / rotate)
\begin{verbentry}{
  style = {default},
  toc = {drehen (to turn / rotate)},
  name = {drehen},
  english = {to turn / rotate},
  type = {regular},
  auxiliary = {haben}
}


\vspace{5pt}
\hrule
\vspace{5pt}

\begin{tabular}{rl}
\multicolumn{2}{l}{\textbf{PRÄSENS}}\\
ich & \textbf{drehe}\\
du & \textbf{drehst}\\
er/sie/es & \textbf{dreht}\\
wir & \textbf{drehen}\\
ihr & \textbf{dreht}\\
sie/Sie & \textbf{drehen}\\
\end{tabular}
\hfill
\begin{tabular}{rl}
\multicolumn{2}{l}{\textbf{PRÄTERITUM}}\\
ich & \textbf{drehte}\\
du & \textbf{drehtest}\\
er/sie/es & \textbf{drehte}\\
wir & \textbf{drehten}\\
ihr & \textbf{drehtet}\\
sie/Sie & \textbf{drehten}\\
\end{tabular}
\hfill
\begin{tabular}{rl}
\multicolumn{2}{l}{\textbf{IMPERATIV}}\\
& \textbf{dreh(e) (du)!}\\
& \textbf{drehen wir!}\\
& \textbf{dreht (ihr)!}\\
& \textbf{drehen Sie!}\\
\end{tabular}

\vspace{10pt}

\begin{tabular}{rl}
\multicolumn{2}{l}{\textbf{PERFEKT}}\\
& haben + \textbf{gedreht}\\
\end{tabular}
\hfill
\begin{tabular}{rl}
\multicolumn{2}{l}{\textbf{FUTURE I}}\\
& werden + \textbf{drehen}\\
\end{tabular}
\hfill
\begin{tabular}{rl}
\multicolumn{2}{l}{\textbf{PLUSQUAMPERFEKT}}\\
& hatten + \textbf{gedreht}\\
\end{tabular}

\vspace{10pt}

\begin{tabular}{lll}
\multicolumn{3}{l}{\textbf{MIT PRÄPOSITIONEN}}\\
& \textbf{drehen an} + Dat & (to turn / adjust)\\
& \textbf{drehen um} + Akk & (to turn around)\\
& \textbf{drehen nach} + Dat & (to turn toward)\\
\end{tabular}
\hfill
\begin{tabular}{lll}
\multicolumn{3}{l}{\textbf{TRENNBARE VERBEN}}\\
& \textbf{ab$|$drehen} & (to turn off)\\
& \textbf{auf$|$drehen} & (to turn up)\\
& \textbf{um$|$drehen} & (to turn around)\\
& \textbf{zu$|$drehen} & (to close / screw shut)\\
\end{tabular}
\vspace{-5pt}
\end{verbentry}

% duzen (to address someone with ``du")
\begin{verbentry}{
  style = {default},
  toc = {duzen (to address someone with ``du")},
  name = {duzen},
  english = {to address someone with \enquote{du}},
  type = {regular},
  auxiliary = {haben}
}


\vspace{5pt}
\hrule
\vspace{5pt}

\begin{tabular}{rl}
\multicolumn{2}{l}{\textbf{PRÄSENS}}\\
ich & \textbf{duze}\\
du & \textbf{duzt}\\
er/sie/es & \textbf{duzt}\\
wir & \textbf{duzen}\\
ihr & \textbf{duzt}\\
sie/Sie & \textbf{duzen}\\
\end{tabular}
\hfill
\begin{tabular}{rl}
\multicolumn{2}{l}{\textbf{PRÄTERITUM}}\\
ich & \textbf{duzte}\\
du & \textbf{duztest}\\
er/sie/es & \textbf{duzte}\\
wir & \textbf{duzten}\\
ihr & \textbf{duztet}\\
sie/Sie & \textbf{duzten}\\
\end{tabular}
\hfill
\begin{tabular}{rl}
\multicolumn{2}{l}{\textbf{IMPERATIV}}\\
& \textbf{duze (du)!}\\
& \textbf{duzen wir!}\\
& \textbf{duzt (ihr)!}\\
& \textbf{duzen Sie!}\\
\end{tabular}

\vspace{10pt}

\begin{tabular}{rl}
\multicolumn{2}{l}{\textbf{PERFEKT}}\\
& haben + \textbf{geduzt}\\
\end{tabular}
\hfill
\begin{tabular}{rl}
\multicolumn{2}{l}{\textbf{FUTURE I}}\\
& werden + \textbf{duzen}\\
\end{tabular}
\hfill
\begin{tabular}{rl}
\multicolumn{2}{l}{\textbf{PLUSQUAMPERFEKT}}\\
& hatten + \textbf{geduzt}\\
\end{tabular}

\vspace{10pt}

\begin{tabular}{lll}
\multicolumn{3}{l}{\textbf{MIT PRÄPOSITIONEN}}\\
& \textbf{duzen mit} + Dat & (to use "du" with someone)\\
\end{tabular}
\hfill
\begin{tabular}{lll}
\multicolumn{3}{l}{\textbf{TRENNBARE VERBEN}}\\
& \textbf{---} & \\
\end{tabular}
\vspace{-5pt}
\end{verbentry}
